\section{Nucleo del banco: dinero, ledger e invariante}

El paquete \codigo{core} concentra el dominio del mini-banco: como se representa
el dinero, donde viven los saldos, como se mueve dinero de una cuenta a otra sin
condiciones de carrera y que invariante debe respetarse siempre. Es el corazon
de la rubrica de Consistencia (25 puntos): el saldo global debe conservarse, sin
inconsistencias, aun bajo carga concurrente. Todo lo que sigue se apoya en una
decision de diseno simple pero estricta: el dinero nunca se representa con punto
flotante.

\subsection{Dinero exacto en centavos}

Un saldo jamas se guarda como \codigo{double}: se guarda como un entero de
centavos (\codigo{long}). Asi se evitan los errores de redondeo binario tipicos
del punto flotante, donde una cantidad como 0.10 no tiene representacion exacta y
las sumas acumulan deriva. Al trabajar en centavos enteros, toda operacion
aritmetica (sumar, restar, comparar) es exacta y reproducible en cualquier nodo.

Por eso los saldos viven como \codigo{long} en \codigo{core/Account.java}
(campo \codigo{balanceCents}) y los montos de cada transferencia tambien, en el
campo \codigo{cents} de \codigo{core/Transfer.java}. La clase
\codigo{core/Money.java} no almacena dinero ni hace aritmetica de saldos: es un
ayudante estatico (constructor privado, sin instancias) que centraliza la unica
conversion entre la cadena decimal que ve el humano (con dos decimales) y la
representacion interna en centavos. Se extrajo a proposito para que los
endpoints, el cargador de datos y el codec del cluster nunca incrusten
aritmetica de \codigo{BigDecimal} ad hoc.

La conversion de texto a centavos usa \codigo{BigDecimal} con redondeo
\codigo{HALF\_UP}, de modo que el parseo nunca introduce error de coma flotante:

\begin{lstlisting}[language=Java,caption={Conversion exacta de decimal a centavos en Money.toCents}]
public static long toCents(String decimal) {
    return new BigDecimal(decimal)
            .movePointRight(2)
            .setScale(0, RoundingMode.HALF_UP)
            .longValueExact();
}
\end{lstlisting}

El camino inverso, \codigo{Money.toDecimal(long)}, reconstruye la cadena con
exactamente dos decimales usando \codigo{movePointLeft(2)} y
\codigo{setScale(2)}. La aritmetica real de mover dinero (la resta del origen y
la suma al destino) no ocurre en \codigo{Money}, sino en el ledger, sobre esos
mismos \codigo{long} de centavos.

\subsection{Cuentas y ledger}

Una \codigo{core/Account.java} es una cuenta del banco: su id y su titular
(\codigo{owner}) son inmutables, y el unico estado mutable es el saldo en
centavos. Ese saldo solo puede leerse o escribirse manteniendo el monitor
intrinseco de la propia cuenta; por eso tanto el lector \codigo{balanceCents()}
como el escritor \codigo{setBalanceCents(long)} estan declarados
\codigo{synchronized}. No hay un objeto candado por campo: el ledger toma los
monitores intrinsecos de las dos cuentas que participan en una transferencia,
ordenados por id, para mantener el protocolo libre de deadlock.

El \codigo{core/Ledger.java} guarda el conjunto de cuentas y realiza la mecanica
cruda de mover dinero entre dos de ellas; deliberadamente no sabe nada de
secuencias, log de commits ni replicacion (eso vive en \codigo{Bank}). Las
cuentas se almacenan en un \codigo{ConcurrentHashMap<Integer, Account>}, de modo
que las lecturas (\codigo{get}) y el alta de cuentas (\codigo{add}) son libres de
candado. La consulta de un saldo individual es simplemente
\codigo{ledger.get(id).balanceCents()}. La suma global se obtiene con
\codigo{totalCents()}, que recorre todas las cuentas acumulando sus saldos como
una lectura consistente: aun bajo movimientos concurrentes ese total es
exactamente el valor cargado al inicio, porque el escaneo toma un snapshot
atomico (lo detalla la subseccion \emph{Snapshot atomico del total}). Observarlo
bajo carga es la forma en que el sistema verifica que las transferencias
conservan el dinero.

Todos los nodos parten del mismo estado gracias a \codigo{core/Dataset.java},
que carga el dataset compartido de cuentas en un ledger. El archivo es un CSV
plano de cuatro columnas sin encabezado
(\codigo{firstName,lastName1,lastName2,balance}) y sin columna de id: el id de
cada cuenta es su posicion en el archivo contando desde \codigo{idBase}
(normalmente 1), asi que la primera linea de datos se parsea, no se omite. El
titular se forma uniendo los tres campos de nombre con espacios, y el saldo se
convierte de su texto decimal a centavos enteros con \codigo{Money.toCents}, de
modo que ningun punto flotante toca jamas un saldo. Como cada nodo lee el mismo
archivo, todos arrancan con saldos identicos, y por eso la comprobacion de
conservacion de dinero a lo largo del cluster es significativa.

\subsection{Transferencia atomica}

La orquestacion de una transferencia solicitada por el cliente vive en
\codigo{Bank.transfer(int, int, long)} de \codigo{core/Bank.java}. El banco
primero aplica las validaciones que son politica del banco (no mecanica del
ledger): rechaza una transferencia a la misma cuenta y un monto no positivo.
Solo si esas comprobaciones pasan, delega el debito y el credito al ledger y, en
el mismo paso atomico (dentro del candado del movimiento), le asigna el siguiente
numero de secuencia; luego la registra en el log y la entrega a los listeners:

\begin{lstlisting}[language=Java,caption={Bank.transfer: politica, movimiento delegado y commit}]
public long transfer(int from, int to, long cents) throws TransferException {
    if (from == to) {
        throw TransferException.selfTransfer();
    }
    if (cents <= 0L) {
        throw TransferException.badAmount();
    }
    // La secuencia se asigna DENTRO del candado del movimiento (atomica con el).
    long seq = ledger.move(from, to, cents, sequence::incrementAndGet);
    Transfer committed = new Transfer(seq, from, to, cents);
    // ...
}
\end{lstlisting}

La parte critica para la consistencia, la atomicidad y la ausencia de deadlock,
no esta en \codigo{Bank} sino en el metodo privado \codigo{apply} del ledger, al
que llega \codigo{ledger.move}. Ese metodo busca ambas cuentas y, para mover el
dinero como un solo paso atomico, toma los dos monitores en un orden fijo
determinado por el id (siempre el menor primero). Ese orden global es lo que
impide el deadlock: si dos transferencias inversas tocan el mismo par de cuentas
en sentidos opuestos, ambas intentan adquirir los candados en la misma
secuencia, de modo que nunca se bloquean mutuamente esperando candados cruzados.

\begin{lstlisting}[language=Java,caption={Ledger.apply: orden de candados por id y movimiento atomico}]
Account first = (from < to) ? src : dst;
Account second = (from < to) ? dst : src;
synchronized (first) {
    synchronized (second) {
        if (enforceFunds && src.balanceCents() < cents) {
            throw TransferException.lowBalance(from);
        }
        src.setBalanceCents(src.balanceCents() - cents);
        dst.setBalanceCents(dst.balanceCents() + cents);
        // Estampa la secuencia aqui dentro: el orden de commit coincide
        // con el orden en que se movio el dinero.
        return (sequencer != null) ? sequencer.getAsLong() : 0L;
    }
}
\end{lstlisting}

Como ambos monitores se mantienen durante todo el movimiento, el par de saldos
cambia junto, sin que otro hilo observe un estado intermedio donde el dinero ya
salio del origen pero aun no entro al destino. El rechazo nunca se devuelve como
valor de retorno, sino como una \codigo{core/TransferException.java}, una
excepcion verificada con un \codigo{code()} de un vocabulario pequeno y cerrado:
\codigo{self\_transfer}, \codigo{bad\_amount}, \codigo{no\_such\_account} y
\codigo{low\_balance}. Asi un llamador no puede continuar por accidente despues
de un movimiento rechazado, y la capa HTTP mapea esos codigos a estados y a un
cuerpo JSON de error. La \codigo{core/Transfer.java} resultante es un
\codigo{record} inmutable (\codigo{seq}, \codigo{from}, \codigo{to},
\codigo{cents}): es la unidad de replicacion y durabilidad que el codec del
cluster y el journal en GCS serializan.

\subsection{Snapshot atomico del total (lectura consistente)}

Sumar las 820k cuentas mientras el banco recibe transferencias no es inocuo: un
escaneo ingenuo puede caer justo entre el debito y el credito de una misma
transferencia (origen ya restado, destino aun no sumado) y reportar un total
\emph{roto} que no corresponde a ningun estado real. Para que
\codigo{Ledger.totalCents()} sea siempre una lectura consistente, el ledger anade
un segundo mecanismo de sincronizacion, separado de los monitores de cuenta: un
\codigo{ReentrantReadWriteLock} llamado \codigo{snapshotLock}.

La clave es una inversion deliberada de los lados del candado. Cada movimiento de
dinero (los caminos \codigo{move}, \codigo{settle} y el \codigo{apply} privado
que comparten) toma el lado \textbf{COMPARTIDO} (read) de \codigo{snapshotLock}
mientras hace su debito y su credito; el escaneo de \codigo{totalCents()} toma el
lado \textbf{EXCLUSIVO} (write). Es decir, los lectores-del-candado son los
escritores-del-dinero, y el escritor-del-candado es el unico lector global. Como
el lado compartido admite cualquier numero de titulares simultaneos, todas las
transferencias siguen ejecutandose en paralelo entre si; pero el lado exclusivo
excluye a todas, de modo que la suma no puede solaparse jamas con una
transferencia a medio aplicar. El escaneo observa, por construccion, un estado
donde cada transferencia esta completa o aun no empezada: un \emph{snapshot}
atomico.

\begin{lstlisting}[language=Java,caption={Ledger: lado compartido en el movimiento, lado exclusivo en el escaneo}]
// En apply(...): el debito+credito corre bajo el lado COMPARTIDO,
// asi muchos transferencias avanzan a la vez. El readLock se toma ANTES
// de los monitores de cuenta (el monitor de cuenta es el sumidero del grafo).
snapshotLock.readLock().lock();
try {
    synchronized (first) {
        synchronized (second) {
            src.setBalanceCents(src.balanceCents() - cents);
            dst.setBalanceCents(dst.balanceCents() + cents);
            return (sequencer != null) ? sequencer.getAsLong() : 0L;
        }
    }
} finally {
    snapshotLock.readLock().unlock();
}

// En totalCents(): el escaneo corre bajo el lado EXCLUSIVO, que excluye
// a todo movimiento, por lo que nunca ve un debito sin su credito.
public long totalCents() {
    snapshotLock.writeLock().lock();
    try {
        long sum = 0L;
        for (Account a : accounts.values()) {
            sum += a.balanceCents();
        }
        return sum;
    } finally {
        snapshotLock.writeLock().unlock();
    }
}
\end{lstlisting}

Las dos capas de candados componen sin riesgo de deadlock porque forman un grafo
dirigido aciclico: el \codigo{readLock} se adquiere \emph{antes} que los monitores
de cuenta, y el monitor de cuenta es siempre el sumidero (nunca se toma el
\codigo{snapshotLock} teniendo ya un monitor de cuenta). Asi, el orden por id que
ya evitaba el deadlock entre transferencias inversas se mantiene intacto, y el
nuevo candado solo agrega un nivel por encima.

\textbf{Por que el invariante queda exacto en los tres nodos.} Como toda
transferencia es suma-cero (un debito acompanado de un credito igual), el total
verdadero es invariante: no cambia jamas. Y como el escaneo de cada nodo corre
bajo el lado exclusivo, no puede capturar una transferencia a medio aplicar; por
tanto reporta exactamente ese total invariante, no un valor transitorio. Los tres
nodos parten del mismo dataset, de modo que los tres escaneos atomicos devuelven
la misma constante aun bajo carga. Eso es justo lo que el sello de consistencia
del dashboard compara: con el snapshot atomico, el sello permanece verde bajo
trafico sostenido en lugar de parpadear por lecturas rotas espurias.

\textbf{Por que \codigo{fair} evita la inanicion del escaneo.} El
\codigo{snapshotLock} se construye en modo justo (\codigo{new
ReentrantReadWriteLock(true)}). Sin justicia, un flujo continuo de transferencias
(lectores-del-candado) podria mantener el lado compartido siempre tomado y
\emph{matar de hambre} al escaneo (escritor-del-candado), que nunca conseguiria
el lado exclusivo. Con justicia, el escritor que espera se atiende en orden de
llegada, asi que el endpoint \codigo{/panel} no puede colgarse bajo carga
sostenida de transferencias.

\subsection{Invariante de saldo y numero de secuencia}

El invariante central es la conservacion del dinero: como cada movimiento es un
debito en una cuenta acompanado de un credito por la misma cantidad en otra, la
suma global de saldos permanece constante. Ese es precisamente el valor que
expone \codigo{Ledger.totalCents()}, leido como un snapshot atomico, y su
constancia bajo carga es lo que demuestra que las transferencias no crean ni
destruyen dinero. El ledger se
mantiene a proposito estrecho (solo aplica debito y credito equilibrados) para
que ese invariante sea facil de razonar.

El numero de secuencia es la otra mitad de la consistencia, y no vive en el
ledger sino en \codigo{core/Bank.java}, como un \codigo{AtomicLong sequence}. Es
monotono y solo avanza en un commit real: la secuencia
(\codigo{sequence.incrementAndGet()}) se asigna \emph{dentro} del candado del
movimiento, atomica con el debito y el credito, de modo que una transferencia
rechazada nunca consume un numero y, ademas, el orden de las secuencias coincide
con el orden en que el dinero se movio. Eso importa para las replicas: como sirven
lecturas mientras reproducen en orden de secuencia, estampar la secuencia junto al
movimiento evita que una cuenta se vea \emph{transitoriamente} en negativo (un
debito reproducido antes del credito que lo fondea). Esa secuencia, estampada en
cada \codigo{Transfer}, es la base de la replicacion y de la verificacion de
consistencia: ordena los commits y permite a una replica saber hasta donde esta al
dia. En el camino de replica,
\codigo{Bank.applyReplicated(Transfer)} no genera secuencia nueva sino que la
recibe con la entrada y eleva los contadores con
\codigo{accumulateAndGet(t.seq(), Math::max)}, de modo que reaplicar una entrada
duplicada nunca retrocede el avance.

Cada commit aceptado se registra en el \codigo{core/TransferLog.java}, un
historial append-only y thread-safe. Conserva la secuencia ordenada de commits
para que la alimentacion de replicas pueda reproducir el historial de catch-up a
una replica recien conectada. Las altas ocurren en el camino de commit a la tasa
plena de transferencias (20\% del trafico del concurso), asi que la lista interna
es un \codigo{ArrayList} protegido por el monitor del objeto: anexar es O(1)
amortizado. Una lista copy-on-write seria un error aqui porque copiaria todo el
arreglo en cada anexion, volviendo la carga sostenida un costo de copia
cuadratico. Los escaneos de catch-up (\codigo{since}) son raros (solo cuando una
replica se conecta) y devuelven una copia fresca de la cola que coincide, de modo
que el llamador la itera sin sostener el candado.

Por ultimo, \codigo{core/CommitListener.java} es la interfaz funcional que
notifica un commit. Los listeners se registran con \codigo{Bank.addCommitListener}
y, en el codigo, el fan-out ocurre unicamente en el camino de cliente del lider:
al final de \codigo{Bank.transfer} se recorre la lista (un
\codigo{CopyOnWriteArrayList}, para no bloquear el registro) y se invoca
\codigo{onCommit(committed)} en cada uno, en orden de registro. Los listeners
concretos son el journal en GCS (durabilidad) y la alimentacion de replicas
(replicacion en vivo). La via \codigo{applyReplicated} no vuelve a notificar a
los listeners, porque ya esta consumiendo el flujo del lider y no debe
re-emitirlo.
